\section{Introduction}
\label{sec:intro}

Interactive AI agents are increasingly used to operate remote systems through terminal commands. Large Language Model (LLM)-based agents execute frequent chains of commands in which each action depends on the parsed output of the preceding one \cite{Liu2025}. Because of their reliance on the action-observation loop, agents are highly sensitive to transport-layer latency. Elevated latency can trigger timeouts or delayed observations, causing agents to misunderstand slow responses as execution failures and leading to errors \cite{Xie2024}. As a result, the responsiveness and robustness of the underlying remote control protocol directly affect the effectiveness of the agent.

Today, SSHv2 remains the de facto protocol for remote shell access and command execution. However, its TCP-based transport can become a bottleneck under network impairment. Head-of-line (HOL) blocking can lead to terminal freezes or delayed command feedback \cite{Scharf2007}. In addition, TCP connections are coupled to network addresses, making SSHv2 sessions vulnerable to disconnection when the client changes networks, such as when moving from Wi-Fi to cellular connectivity. These limitations are problematic for interactive AI agents, which require continuous feedback from remote environments to make timely decisions. In response, alternative remote control protocols have emerged. Mosh replaces the traditional byte stream model with a UDP-based design and terminal state synchronization. It aims to improve responsiveness under packet loss and support roaming across networks \cite{Winstein2012}. In contrast, SSH3 explores a redesign of secure remote control over HTTP/3 and QUIC \cite{michel2023}. This architectural shift can reduce connection setup latency while enabling stream multiplexing and connection migration.

However, these design advantages do not imply that a single protocol will dominate all interactive AI agent interactions. Remote control performance for interactive AI agents is workload-dependent. Existing evaluations of remote terminal systems typically focus on human users, bulk data transfer, or general terminal responsiveness \cite{Xie2024}. To the best of our knowledge, they do not evaluate SSHv2, Mosh, and SSH3 under AI agent interaction patterns. This paper investigates four research questions. (RQ1) Which protocol provides the lowest interactive latency for interactive AI agent command execution? (RQ2) Which protocol is most robust under delay, jitter, and packet loss? (RQ3) Does QUIC-based SSH3 improve remote control performance compared with TCP-based SSHv2? (RQ4) Which protocol is best suited for each class of interactive AI agent workloads?

To answer these questions, we evaluate SSHv2, Mosh, and SSH3 using a controlled testbed composed of an interactive AI agent client, a tc/netem-based network emulator, and a remote server. The evaluation covers four classes of workloads: command loops, continuous monitoring, interactive editing, and large-output commands. Experiments are conducted under low, medium, high impairment levels, plus a real-world Internet condition. Our results show that protocol performance is strongly workload-dependent. SSH3 achieves the fastest session establishment and lowest latency for lightweight commands, but its userspace QUIC implementation underperforms kernel TCP for bulk-output delivery. Mosh achieves the lowest apparent latency across all workloads, but delivers only 2--6\% of expected output for large commands, rendering it unsuitable for AI agents that must parse complete responses. SSHv2 remains the most reliable choice for commands producing large output. 

This paper makes the following contributions:

\begin{itemize}
    \item We present the first comparative evaluation of SSHv2, Mosh, and SSH3 for interactive AI agent remote control workloads.
    \item We design a workload suite that captures common AI agent interaction patterns, including command loops, continuous monitoring, interactive editing, and large-output commands.
    \item We quantify per-command latency, latency variance, and tail behavior of each protocol under controlled network impairment and a real-world Internet condition.
    \item We derive practical guidelines for selecting remote control protocols in interactive AI agent deployments.
\end{itemize}

The remainder of this paper is organized as follows. Section~\ref{sec:background} presents the background and related work. Section~\ref{sec:proposal} describes the measurement methodology, workloads, network scenarios, and metrics. Section~\ref{sec:evaluation} reports and analyzes the experimental results. Finally, Section~\ref{sec:conclusion} concludes the paper.